\subsection{The group of feasible edge replacements}\label{subsec:fer}

I start this subsection by setting up a combinatorial puzzle.
A graph is given, and one may remove an edge and add it wherever there is no edge (trivially, one can add it back to its original position), but the new graph must be isomorphic to the starting graph.
Such an action is called a \emph{feasible edge replacement}, which is formalized after these examples.
Can one attain every possible configuration of the graph by repeatedly changing edges like this?
It is clear that the answer depends on the structure of the particular starting graph.
For instance, a cycle of four edges has no valid non-trivial movements since, having removed an edge, the only possibility that preserves the cycle is putting the edge back from where it was taken.
Graphically, Figure~\ref{fig:4_cycles} depicts two isomorphic graphs that cannot be transformed into each other by feasible edge replacements.

\begin{figure}[htbp]
\centering
% Left Subfigure: "U-shape" path
\begin{subfigure}[b]{0.4\textwidth}
    \centering
    \begin{tikzpicture}[scale=2]
        % Define coordinates for a square
        \node[vertex] (TL) at (0,1) {};
        \node[vertex] (TR) at (1,1) {};
        \node[vertex] (BL) at (0,0) {};
        \node[vertex] (BR) at (1,0) {};

        % Draw the path: Left, Top, Right
        \draw[thick] (BL) -- (TL) -- (TR) -- (BR) -- (BL);
    \end{tikzpicture}
\end{subfigure}
\hfill
% Right Subfigure: "Z-shape" / Diagonals
\begin{subfigure}[b]{0.4\textwidth}
    \centering
    \begin{tikzpicture}[scale=2]
        % Define coordinates for a square
        \node[vertex] (TL) at (0,1) {};
        \node[vertex] (TR) at (1,1) {};
        \node[vertex] (BL) at (0,0) {};
        \node[vertex] (BR) at (1,0) {};

        % Draw the path: Bottom-Left to Top-Left, Top-Left to Bottom-Right, Bottom-Right to Top-Right
        % Based on your description: "top edge and the two diagonals"
        \draw[thick] (BL) -- (TR); % Diagonal 1
        \draw[thick] (TR) -- (TL); % Top edge
        \draw[thick] (TL) -- (BR); % Diagonal 2
        \draw[thick] (BL) -- (BR);
    \end{tikzpicture}
\end{subfigure}
\caption{Two $4$-cycles on the same vertex set.}
\label{fig:4_cycles}
\end{figure}

Now for the puzzle: the two paths in Figure~\ref{fig:4_paths} are isomorphic.
Is it possible to transform one into the other by a sequence of feasible edge replacements?

\begin{figure}[htbp]
\centering
% Left Subfigure: "U-shape" path
\begin{subfigure}[b]{0.4\textwidth}
    \centering
    \begin{tikzpicture}[scale=2]
        % Define coordinates for a square
        \node[vertex] (TL) at (0,1) {};
        \node[vertex] (TR) at (1,1) {};
        \node[vertex] (BL) at (0,0) {};
        \node[vertex] (BR) at (1,0) {};

        % Draw the path: Left, Top, Right
        \draw[thick] (BL) -- (TL) -- (TR) -- (BR);
    \end{tikzpicture}
\end{subfigure}
\hfill
% Right Subfigure: "Z-shape" / Diagonals
\begin{subfigure}[b]{0.4\textwidth}
    \centering
    \begin{tikzpicture}[scale=2]
        % Define coordinates for a square
        \node[vertex] (TL) at (0,1) {};
        \node[vertex] (TR) at (1,1) {};
        \node[vertex] (BL) at (0,0) {};
        \node[vertex] (BR) at (1,0) {};

        % Draw the path: Bottom-Left to Top-Left, Top-Left to Bottom-Right, Bottom-Right to Top-Right
        % Based on your description: "top edge and the two diagonals"
        \draw[thick] (BL) -- (TR); % Diagonal 1
        \draw[thick] (TR) -- (TL); % Top edge
        \draw[thick] (TL) -- (BR); % Diagonal 2
    \end{tikzpicture}
\end{subfigure}
\caption{Two paths of length $3$ on the same vertex set.}
\label{fig:4_paths}
\end{figure}

\begin{definition}[Labeled graph]\label{def:labeled-graph}
    For a set $X$, whose elements are called \emph{labels}, denote by $S_X$ the set of all bijections, called \emph{label permutations}, $X\to X$.
    Under the operation of composition, $S_X$ forms a group called the \emph{symmetric group on $X$}.
    As usual, we use the notation $\rho\cdot\sigma:=\rho\circ\sigma$ for the group operation.
    Endow this set with the topology inherited from the Tychonoff product $X^X$, with $X$ taken as discrete.
    Let $G$ be a graph defined on a vertex set $V$ and let $\lambda\colon V\to X$ be a bijection, called a \emph{labeling} of $G$.
    For the sake of simplicity, I will always write vertices as $v_\ell\in V$ where $\lambda(v_\ell)=\ell\in X$ is the label.

    The pair $(G,\lambda)$ is called a \emph{labeled graph}, and the parameters $V,E$ and $X$ are considered to be implicit.
\end{definition}

Every label permutation $\sigma\in S_X$ induces the graph isomorphism $\varphi_\sigma:=\lambda^{-1}\circ\sigma\circ\lambda\colon V\to V$ of $G$, and every graph isomorphism of $G$ of the form $\varphi\colon V\to V$ (that is, on the same vertex set) induces a label permutation $\sigma_\varphi:=\lambda\circ\varphi\circ\lambda^{-1}\in S_X$.
For this reason, label permutations and isomorphisms are often used interchangeably, but I do not follow this convention for the sake of clarity.

\begin{center}
\begin{tikzpicture}[node distance=1cm and 4cm, auto, >=Stealth]
    % Top Row (Labels X)
    \node (X1) {$X$};
    \node (X2) [right=of X1] {$X$};
    
    % Bottom Row (Vertices V)
    \node (V1) [below=of X1] {$V$};
    \node (V2) [below=of X2] {$V$};

    % Horizontal Arrows
    \draw[->] (X1) -- node {$\sigma$} (X2);
    \draw[->] (V1) -- node [swap] {$\varphi_\sigma$} (V2);

    % Vertical Arrows (Connecting the lines via lambda)
    \draw[->] (V1) -- node {$\lambda$} (X1);
    \draw[->] (X2) -- node {$\lambda^{-1}$} (V2);
\end{tikzpicture}
\end{center}

\begin{lemma}\label{lemma:varphi-iso-iff-image}
    Let $G$ and $G'$ be graphs on the same vertex set $V$.
    Given a bijection $\varphi\colon V\to V$, $\varphi$ is an isomorphism between $G$ and $G'$ if and only if $$\varphi[E(G)]=E(G').$$
\end{lemma}

\begin{proof}
    This is immediate from the definition of graph isomorphism: $\varphi$ is an isomorphism if and only if $v_iv_j\in E(G)\Leftrightarrow \varphi(v_i)\varphi(v_j)\in E(G')$, which is precisely the statement $\varphi[E(G)]=E(G')$.
\end{proof}

We adapt the notation from \cite{Caro2020GraphsIU} and all subsequent papers.
Given a labeled graph $(G,\lambda)$, the \emph{edge-label set of $G$} is defined as
\[
    L_G:=\{ij : v_iv_j\in E(G)\},
\]
where $ij$ and $ji$ are not distinguished.
Thus, $L_G$ is the set of labels induced on the edges of $G$ by $\lambda$.
For each $\sigma\in S_X$, define $G_\sigma$ to be the graph on the same vertex set $V$ with edge set
\[
    E(G_\sigma):=\{v_iv_j : \sigma(i)\sigma(j)\in L_G\}.
\]
In words, $G_\sigma$ is the unique graph on $V$ satisfying $L_{G_\sigma}=\sigma^{-1}[L_G]$, where $\sigma^{-1}[L_G]:=\{\sigma^{-1}(i)\sigma^{-1}(j):ij\in L_G\}$.
In particular, $L_{G_\sigma}=L_G$ for every $\sigma\in S_X$ (the set of edge labels is invariant across all copies), and $G_{\text{id}}=G$.

The relationship between $G_\sigma$ and $\varphi_\sigma$ is as follows.
By Lemma~\ref{lemma:varphi-iso-iff-image}, $\varphi_\sigma$ is an isomorphism $G\to G'$ if and only if $\varphi_\sigma[E(G)]=E(G')$.
Since $\varphi_\sigma(v_i)=v_{\sigma^{-1}(i)}$, this image equals $\{v_{\sigma^{-1}(i)}v_{\sigma^{-1}(j)}:v_iv_j\in E(G)\}$, whose edge-label set is exactly $\sigma^{-1}[L_G]=L_{G_\sigma}$.
In summary:
\begin{equation}\label{eq:gsigma-varphi}
    \varphi_\sigma \text{ is an isomorphism } G\to G' \quad\Longleftrightarrow\quad G_\sigma = G' \quad\Longleftrightarrow\quad \sigma^{-1}[L_G]=L_{G'}.
\end{equation}

\begin{center}
\begin{tikzpicture}[node distance=1.2cm and 5cm, auto, >=Stealth]
    % Top Row (Label level)
    \node (L1) {$L_G$};
    \node (L2) [right=of L1] {$L_{G_\sigma}=\sigma^{-1}[L_G]$};
    
    % Bottom Row (Graph level)
    \node (G1) [below=of L1] {$G$};
    \node (G2) [below=of L2] {$G_\sigma$};

    % Horizontal Arrows
    \draw[|->] (L1) -- node {$\sigma^{-1}$} (L2);
    \draw[->] (G1) -- node [swap] {$\varphi_\sigma$} (G2);

    % Vertical Arrows
    \draw[<->] (G1) -- node {$\lambda$} (L1);
    \draw[<->] (G2) -- node {$\lambda$} (L2);
\end{tikzpicture}
\end{center}

Suppose that $G^c$ denotes the graph complement of $G$ (that is, the graph obtained from switching every adjacency and every non-adjacency in $G$) and let $(e,f)\in E(G)\times E(G^c)$, and $\sigma\in S_X$.
I use the phrase ``$\sigma::(e\to f)$ is a \emph{feasible edge replacement}'' to abbreviate the fact that $\varphi_\sigma$ is a graph isomorphism between $G$ and $G-e+f$, or equivalently, that $\sigma^{-1}[L_G]=L_{G-e+f}$.\footnote{I delay defining the object $\sigma::(e\to f)$ (Definition~\ref{def:fer-object}), for a few paragraphs, purely for expository purposes, as this requires an additional level of abstraction that would otherwise lack motivation.}
If $\varphi_\sigma\in\Aut(G)$, I employ the symbol $\sigma::(\emptyset\to\emptyset)$ to signal that no edge is being moved.
Denote by $\Fer(G)$ the group generated by all $\sigma$ for which there are $e$ and $f$ such that $\sigma::(e\to f)$ is a feasible edge replacement.
These label permutations $\sigma$ are called the \emph{canonical generators} of $\Fer(G)$, and they include permutations that induce automorphisms on $G$ (which can be visualized as an action that leaves the graph looking identical) and permutations that correspond to moving exactly one edge.
Thus, intuitively, each element of the group $\Fer(G)$ corresponds to a ``sequence'' of feasible edge replacements.

Let $$A(G):=\{\sigma\in S_X:\varphi_\sigma\in\Aut(G)\}\simeq\Aut(G)$$ be the subgroup of label permutations that correspond to the automorphisms of $G$.
Hence, the following subgroup structure holds.
$$1\leq A(G)\leq\Fer(G)\leq S_X,$$ where, as usual, the $1$ on the left represents the trivial subgroup of $S_X$ containing just the identity on $X$.

It is often convenient to represent a sequence of feasible edge replacements (that is, an element of $\Fer(G)$) as a product of canonical generators (which, as standard in Algebra, is written right-to-left) paired with a ``list of instructions'' explaining what edges to replace and in what order (which, as standard in English, is written left-to-right).
This requires defining some more symbols.
Before doing this, there is a technical caveat: the labels of the graph's vertices change after each canonical generator gets applied (the labels get permuted); but having the list of edge replacements in the graph's original labeling is desirable if, for example, one wishes to execute a solution to a puzzle by moving edges without having to keep track of what the corresponding permutations are (not to mention the extra effort needed to figure out what those permutations are in the first place).
To circumvent this, I give the following abstract definition (by ``abstract,'' I mean that the operation does not actually require knowledge of the graph $G$ on which it is defined).

\begin{definition}[Edge replacements as objects]\label{def:fer-object}
    Suppose that $\sigma\colon X\to X$ is label permutation and that $x,y,z,w\in X$, and possibly adding subscripts, are labels.
    An \emph{edge replacement} is a tuple of the form $(\sigma, (x, y, z, w))$ and we always use the notation $\sigma::(xy\to zw)$ to represent it.

    If $\{x,y\}=\{z,w\}$, I use the notation $\sigma::(\emptyset\to\emptyset)$ to emphasize that no edge has been moved.

    A \emph{sequence of edge replacements} is a tuple $(\sigma,(x_i,y_i,z_i,w_i)_{i\leq n})$ where $\sigma$ is again a label permutation, and $\{x_i,y_i,z_i,w_i:i\leq n\}\subseteq X$.
    I employ the symbol $$\sigma::(x_0y_0\to z_0w_0)(x_1y_1\to z_1w_1)\cdots(x_ny_n\to z_nw_n)$$ to refer to this object.
\end{definition}

Intuitively, $\sigma::(xy\to zw)$ refers to applying the permutation $\sigma$ to the labels of the vertices and moving the edge $v_xv_y$ to the position of the non-edge $v_zv_w$.\footnote{For readers previously familiar with this notion: the notation $\sigma::(xy\to zw)$ denotes that $\sigma\in\Fer_G(xy\to zw)$ in the sense of \cite{Caro2020GraphsIU}.}

\begin{definition}[Concatenating edge replacements]\label{def:concatenation-fer}
    Now suppose that \begin{align*}
    f_0=\sigma&::(x_0y_0\to z_0w_0)(x_1y_1\to z_1w_1)\cdots(x_ny_n\to z_nw_n)\qquad\text{ and}\\
    f_1=\rho&::(a_0b_0\to c_0d_0)(a_1b_1\to c_1d_1)\cdots(a_mb_m\to c_md_m)
    \end{align*}
    are both sequences of edge replacements.
    Then the \emph{product} $f_0*f_1$, is the sequence of edge replacements \begin{align*}
        f_0*f_1:=\rho\cdot\sigma::(x_0y_0\to z_0w_0)(x_1y_1\to z_1w_1)\cdots\\
        \cdots(x_ny_n\to z_nw_n)(\sigma^{-1}(a_0)\sigma^{-1}(b_0)\to \sigma^{-1}(c_0)\sigma^{-1}(d_0))\cdots\\
        \cdots(\sigma^{-1}(a_m)\sigma^{-1}(b_m)\to\sigma^{-1}(c_m)\sigma^{-1}(d_m)).
    \end{align*}
\end{definition}

In other words, the inverse permutation must be applied to the second sequence of labels to ensure that the final labeling is consistent with the original labeling of the graph.
This relabeling also has a very natural categorical interpretation, in the category of labeled graphs $(G,\lambda)$.
Each edge replacement corresponds to a morphism between labeled graphs, and a sequence of feasible edge replacements corresponds to compositions of such morphisms.

Now, I connect the two notions thus far defined in this subsection and prove that they are indeed consistent with one another.
The following proposition can be inductively generalized to arbitrary finite sequences of edge replacements.

\begin{proposition}\label{prop:fer-is-well-defined}
    Let $(G,\lambda)$ be a labeled graph and suppose that $f_0=\sigma::(xy\to zw)$ and $f_1=\rho::(ab\to cd)$ are feasible edge replacements of $G$.
    Then, their product, $$f_0*f_1=\rho\cdot\sigma::(xy\to zw)(\sigma^{-1}(a)\sigma^{-1}(b)\to\sigma^{-1}(c)\sigma^{-1}(d))$$ is such that the induced map $\varphi_{\rho\cdot\sigma}$ is an isomorphism between $G$ and $$(G-v_xv_y+v_zv_w)-v_{\sigma^{-1}(a)}v_{\sigma^{-1}(b)}+v_{\sigma^{-1}(c)}v_{\sigma^{-1}(d)}.$$
\end{proposition}

Products of feasible edge replacements are called \tfer \emph{objects}, and their label permutations are precisely the elements of $\Fer(G)$.
The preceding proposition shows that $f_0*f_1$ is a \tfer object, which we prove as follows.

\begin{proof}
    Denote $G':=G-v_xv_y+v_zv_w$ and $G'':=G-v_av_b+v_cv_d$.
    At the level of edge labels, this means
    \[
        L_{G'}=(L_G\setminus\{xy\})\cup\{zw\}\qquad\text{ and }\qquad L_{G''}=(L_G\setminus\{ab\})\cup\{cd\}.
    \]

    Since $f_0$ is a feasible edge replacement of $G$, by \eqref{eq:gsigma-varphi}, $G_\sigma=G'$, i.e.,
    \begin{equation}\label{eq:f0-label}
        \sigma^{-1}[L_G]=L_{G'}=(L_G\setminus\{xy\})\cup\{zw\}.
    \end{equation}
    Since $f_1$ is a feasible edge replacement of $G$, again by \eqref{eq:gsigma-varphi}, $G_\rho=G''$, i.e.,
    \begin{equation}\label{eq:f1-label}
        \rho^{-1}[L_G]=L_{G''}=(L_G\setminus\{ab\})\cup\{cd\}.
    \end{equation}
    
    The goal is to compute $(\rho\cdot\sigma)^{-1}[L_G]$.
    Since $(\rho\cdot\sigma)^{-1}=\sigma^{-1}\circ\rho^{-1}$, applying the bijection $\sigma^{-1}$ to both sides of \eqref{eq:f1-label} gives
    \begin{align}
        (\rho\cdot\sigma)^{-1}[L_G]
            &= \sigma^{-1}\bigl[\rho^{-1}[L_G]\bigr] \nonumber\\
            &= \sigma^{-1}\bigl[(L_G\setminus\{ab\})\cup\{cd\}\bigr] \nonumber\\
            &= \bigl(\sigma^{-1}[L_G]\setminus\{\sigma^{-1}(a)\sigma^{-1}(b)\}\bigr)\cup\{\sigma^{-1}(c)\sigma^{-1}(d)\}. \label{eq:product-step}
    \end{align}
    Substituting \eqref{eq:f0-label} into \eqref{eq:product-step} yields
    \[
        (\rho\cdot\sigma)^{-1}[L_G]= \bigl(L_{G'}\setminus\{\sigma^{-1}(a)\sigma^{-1}(b)\}\bigr)\cup\{\sigma^{-1}(c)\sigma^{-1}(d)\} = L_{G'''},
    \]
    where $G''':=G'-v_{\sigma^{-1}(a)}v_{\sigma^{-1}(b)}+v_{\sigma^{-1}(c)}v_{\sigma^{-1}(d)}$.
    By \eqref{eq:gsigma-varphi}, this means $G_{\rho\cdot\sigma}=G'''$, i.e., $\varphi_{\rho\cdot\sigma}$ is an isomorphism $G\to G'''$.
\end{proof}

With the multiplication I defined in Definition~\ref{def:concatenation-fer}, \tfer objects form their own group structure, which I call the \tFer group; clearly, it is isomorphic to $\Fer(G)$.
As far as I know, this group is not explicitly defined in the existing body of literature, and doing so has some practical advantages, especially so from an engineering perspective.
